\makeatletter
\@ifundefined{keepfig}{\newenvironment{keepfig}{\par\edef\keepfig@pi{\the\parindent}\noindent
  \begin{minipage}{\linewidth}\parindent=\keepfig@pi\relax}{\end{minipage}\par}}{}
\makeatother
\section{Основная теорема алгебры}

Докажем, наконец, основную теорему алгебры.

\begin{theorem}[5.6]
Поле $\CC$ комплексных чисел алгебраически замкнуто.
\end{theorem}

Строго говоря, это утверждение не только алгебраическое: оно о поле $\R$, построение которого~--- предмет анализа.
Из анализа нам понадобятся ровно два свойства $\R$, и оба следуют из теоремы о промежуточном значении (непрерывная на
отрезке функция принимает все значения между значениями на концах).

\begin{enumerate}[label=(\roman*),leftmargin=*]
\item \emph{Каждое положительное число $a\in\R$ имеет квадратный корень в $\R$.} Функция $x\mapsto x^2$ непрерывна на
отрезке $[0,b]$, где $b=\max(1,a)$, и $0^2=0<a\le b^2$; значит, $x^2=a$ для некоторого $x\in[0,b]$.
\item \emph{Каждый многочлен нечётной степени с вещественными коэффициентами имеет вещественный корень.} Можно считать
многочлен $p(x)=x^{2k+1}+a_{2k}x^{2k}+\dots+a_0$ унитарным. Так как $p(x)\to+\infty$ при $x\to+\infty$ и $p(x)\to-\infty$
при $x\to-\infty$, найдётся $N>0$ с $p(-N)<0<p(N)$, и $p$ обращается в нуль на отрезке $[-N,N]$.
\end{enumerate}

\begin{keepfig}
\begin{center}
\begin{tikzpicture}[>=stealth,scale=0.95]
\begin{scope}
\draw[->] (-0.3,0) -- (2.4,0) node[below] {$x$};
\draw[->] (0,-0.3) -- (0,3.3) node[left] {$y$};
\draw[thick,cold,domain=0:1.75,samples=40] plot (\x,{\x*\x}) node[right] {$y=x^2$};
\draw[cnew,dashed] (0,1.7) node[left,black] {$a$} -- (1.304,1.7) -- (1.304,0) node[below,black] {$\sqrt a$};
\fill[cnew] (1.304,1.7) circle (1.4pt);
\end{scope}
\begin{scope}[xshift=6.3cm,yshift=1.3cm]
\draw[->] (-2.5,0) -- (2.6,0) node[below] {$x$};
\draw[->] (0,-1.2) -- (0,2.0) node[left] {$y$};
\draw[thick,cold,domain=-2.05:1.9,samples=80] plot (\x,{0.5*\x*\x*\x-1.5*\x+0.5}) node[right] {$y=p(x)$};
\draw[dashed] (-2.05,0) -- (-2.05,-0.73) node[below,font=\footnotesize] {$-N$};
\draw[dashed] (1.9,0) -- (1.9,1.08);
\node[below,font=\footnotesize] at (1.9,0) {$N$};
\foreach \r in {0.3473,1.5321,-1.8794}{\fill[cnew] (\r,0) circle (1.8pt);}
\end{scope}
\end{tikzpicture}
\end{center}
\figcap{Два свойства $\R$, следующие из теоремы о промежуточном значении: квадратный корень из $a>0$ (слева) и корень
многочлена нечётной степени, который меняет знак на отрезке $[-N,N]$ (справа; красные точки~--- корни).}
\end{keepfig}

Первое свойство позволяет извлекать квадратные корни уже в $\CC$.

\begin{lemma}[5.6a]
Каждое число $\alpha=a+bi\in\CC$ ($a,b\in\R$) является квадратом в $\CC$: если $r=\sqrt{a^2+b^2}$, то
\[
\alpha=(c+di)^2,\qquad c=\sqrt{\frac{r+a}2},\quad d=\varepsilon\sqrt{\frac{r-a}2},\qquad
\varepsilon=\begin{cases}\phantom{-}1,&b\ge0,\\-1,&b<0.\end{cases}
\]
\end{lemma}

\begin{proof}
Так как $r\ge|a|$, числа $(r\pm a)/2$ неотрицательны, и корни из них существуют по (i). Далее,
\[
c^2-d^2=\frac{r+a}2-\frac{r-a}2=a,\qquad 2cd=2\varepsilon\sqrt{\frac{(r+a)(r-a)}4}=\varepsilon\sqrt{r^2-a^2}=\varepsilon|b|=b,
\]
поэтому $(c+di)^2=(c^2-d^2)+2cd\,i=a+bi$.
\end{proof}

Формулу легко найти. Равенство $a+bi=(c+di)^2=(c^2-d^2)+2cd\,i$ с вещественными $c,d$ означает систему
\[
\begin{cases}c^2-d^2=a,\\ 2cd=b.\end{cases}
\]
Если $c\ne0$, то $d=\frac b{2c}$ и $c^2-\frac{b^2}{4c^2}=a$, т.\,е.\ $4c^4-4ac^2-b^2=0$. Это квадратное уравнение
относительно $c^2$, и
\[
c^2=\frac{4a\pm\sqrt{16a^2+16b^2}}8=\frac{a\pm\sqrt{a^2+b^2}}2 .
\]
Так как $c\ne0$, имеем $c^2>0$, а $a-\sqrt{a^2+b^2}\le0$; поэтому нужно взять знак плюс: $c=\pm\sqrt{(a+r)/2}$. Тогда
\[
d^2=\frac{b^2}{4c^2}=\frac{b^2}{2(a+r)}=\frac{b^2(r-a)}{2(r^2-a^2)}=\frac{r-a}2
\]
(при $b\ne0$, когда $r^2-a^2=b^2\ne0$; при $b=0$ имеем $a=r>0$ и $d=0=\frac{r-a}2$), и знаки $c$ и $d$ надо согласовать так, чтобы знак $cd$ совпадал со знаком $b$ (ведь $2cd=b$). Случай $c=0$ возможен
лишь при $b=0$, $a\le0$, и тогда $\alpha=(i\sqrt{-a})^2$~--- это тоже частный случай формулы леммы. Например,
\[
\sqrt{3+4i}=\pm(2+i),\qquad \sqrt{-5+12i}=\pm(2+3i),\qquad \sqrt i=\pm\frac{1+i}{\sqrt2}.
\]

\begin{proof}[Доказательство теоремы 5.6]
Определим $\CC$ как поле разложения многочлена $X^2+1$ над $\R$ и обозначим через $i$ его корень; тогда
$\CC=\R[i]=\{a+bi\mid a,b\in\R\}$. По \rf{1.45} достаточно доказать, что каждый многочлен $f\in\R[X]$ раскладывается в
$\CC[X]$ на линейные множители.

Пусть $E$~--- поле разложения многочлена $f(X)\,(X^2+1)$ над $\CC$; так как $\CC$ порождается над $\R$ корнем $X^2+1$,
это и поле разложения $f(X)(X^2+1)$ над $\R$. Нужно доказать, что $E=\CC$. Поле $E$ является и полем разложения
многочлена $h$~--- произведения \emph{различных} унитарных неприводимых множителей $f(X)(X^2+1)$ в $\R[X]$ (корни у этих
двух многочленов одни и те же). В характеристике $0$ неприводимые многочлены сепарабельны, а различные унитарные
неприводимые многочлены не имеют общих корней (\rf{2.17}); поэтому $h$ сепарабелен, и $E/\R$~--- расширение Галуа (\rf{3.10}).
Пусть $G=\Gal(E/\R)$ и $H$~--- силовская $2$-подгруппа $G$.

\emph{Шаг 1: $G$~--- $2$-группа.} Пусть $M=E^H$. По основной теореме $[E:M]=|H|$, поэтому $[M:\R]=(G:H)$~--- нечётное
число. Для $\alpha\in M$ степень $[\R[\alpha]:\R]$ делит $[M:\R]$ (\rf{1.20}) и потому нечётна: минимальный многочлен
$\alpha$ над $\R$ имеет нечётную степень. По (ii) у него есть вещественный корень $x_0$, т.\,е.\ он делится на $X-x_0$; а
так как он неприводим, он равен $X-x_0$. Значит, $\alpha=x_0\in\R$. Итак, $M=\R$, и по основной теореме
$H=\Gal(E/E^H)=\Gal(E/\R)=G$.

\emph{Шаг 2: $\Gal(E/\CC)=1$.} Подгруппа $\Gal(E/\CC)\subset G$ имеет порядок $2^m$. Предположим, что $m\ge1$. Конечная
$p$-группа порядка $p^m$ имеет подгруппы всех порядков $p^k$, $0\le k\le m$. (Индукция по $m$: центр $p$-группы
нетривиален по формуле классов, в нём есть элемент $z$ порядка $p$; подгруппа $\langle z\rangle$ нормальна, и прообразы
подгрупп порядка $p^{k-1}$ факторгруппы по $\langle z\rangle$ дают подгруппы порядка $p^k$.) Поэтому в $\Gal(E/\CC)$ есть
подгруппа $N$ индекса $2$. Её неподвижное поле $E^N$ имеет степень $2$ над $\CC$, а квадратичное расширение поля
характеристики $\ne2$ получается присоединением квадратного корня: $E^N=\CC[\sqrt d]$ для некоторого $d\in\CC$ (\rf{3.25}).
Но по лемме \rf{5.6a} $\sqrt d\in\CC$, так что $E^N=\CC$~--- противоречие. Значит, $m=0$, $E=\CC$, и $f$ раскладывается в
$\CC[X]$ на линейные множители.
\end{proof}

\begin{keepfig}
\begin{center}
\begin{tikzpicture}[yscale=0.9]
\node[cold] (E) at (0,3.2) {$E$};
\node[cnew] (M) at (-2.2,1.6) {$M=E^H$};
\node[cnew] (EN) at (2.0,2.2) {$E^N$};
\node[cold] (C) at (2.0,1.0) {$\CC=\R[i]$};
\node[cold] (R) at (0,0) {$\R$};
\draw (R) -- (M);
\node[font=\scriptsize,align=right,anchor=east] at (-1.25,0.55) {нечётная\\ степень $(G:H)$};
\draw (M) -- node[left,font=\scriptsize]{$|H|$} (E);
\draw (R) -- node[below right,font=\scriptsize]{$2$} (C);
\draw[cnew] (C) -- node[right,font=\scriptsize]{$2$} (EN);
\draw (EN) -- (E);
\end{tikzpicture}
\end{center}
\figcap{Шаг 1: поле $M=E^H$ нечётной степени над $\R$ совпадает с $\R$, поэтому $G=H$~--- $2$-группа. Шаг 2: подгруппа
индекса $2$ в $\Gal(E/\CC)$ дала бы квадратичное расширение $E^N$ поля $\CC$, а их нет, так как в $\CC$ извлекаются
квадратные корни.}
\end{keepfig}

\begin{corollary}[5.7]
\begin{enumerate}[label=\textup{(\alph*)}]
\item Поле $\CC$~--- алгебраическое замыкание поля $\R$.
\item Множество всех алгебраических чисел (комплексных чисел, алгебраических над $\Q$) является алгебраическим
замыканием поля $\Q$.
\end{enumerate}
\end{corollary}

\begin{proof}
(a) Поле $\CC$ алгебраически замкнуто (\rf{5.6}) и алгебраично над $\R$, так как $[\CC:\R]=2$; это и есть определение
алгебраического замыкания (\rf{1.44}). (b) По \rf{1.47}, если $\Om$ алгебраически замкнуто и $F\subset\Om$, то элементы $\Om$,
алгебраические над $F$, образуют алгебраическое замыкание $F$; применяем это к $\Om=\CC$ и $F=\Q$.
\end{proof}

\begin{aside}[5.7a]
Есть и короткое аналитическое доказательство теоремы \rf{5.6}. Пусть $f(z)=z^m+a_{m-1}z^{m-1}+\dots+a_0\in\CC[z]$, $m\ge1$, и
$B=\max|a_k|$. При $|z|\ge\max(1,2mB)$ имеем $|a_{m-1}z^{m-1}+\dots+a_0|\le mB|z|^{m-1}\le\frac12|z|^m$, откуда
$|f(z)|\ge\frac12|z|^m\to\infty$ при $|z|\to\infty$. Если бы у $f$ не было корней, функция $1/f$ была бы голоморфной на
всей плоскости и ограниченной (на большом круге~--- по непрерывности, вне его~--- по оценке), а значит, постоянной по
теореме Лиувилля,~--- противоречие. Приведённое выше алгебраическое доказательство~--- упрощение Э.~Артином более ранних
доказательств; из анализа в нём используются лишь свойства (i) и (ii).
\end{aside}
